\subsection{PID 匿名身份机制}

PID（Pseudonymous Identifier）即伪身份标识或匿名身份标识，是隐私保护系统中常用的一种身份解耦机制。其核心思想是：系统内部不直接使用真实用户身份参与通信和认证，而是为用户生成一个与真实身份存在受控映射关系的匿名标识。在日常业务过程中，通信参与方只能看到订单相关信息或匿名身份信息，不能直接获得用户真实身份。

在 FoodShield 系统中，用户注册后由平台服务器生成 PID。PID 由主密钥、用户真实标识、随机数和设备摘要等字段通过 HMAC 机制生成，其抽象形式如下：

\begin{equation}
\mathrm{PID}=\operatorname{HMAC}(K_{\mathrm{master}},\, \mathit{userID} \parallel r)
\end{equation}

其中，$K_{\mathrm{master}}$ 为平台主密钥，$\mathit{userID}$ 为用户真实身份标识，$r$ 为随机数。引入随机数 $r$ 的作用在于实现\textbf{跨订单不可链接性}（Cross-Order Unlinkability）：即使同一用户在不同订单中使用 PID，由于每次生成的随机数不同，攻击者无法通过比较 PID 值来判断两个订单是否属于同一用户。这一性质对于防止跨订单的用户画像构建和身份关联攻击具有重要意义。

PID 的安全性可以从以下两个角度分析：

\begin{itemize}[leftmargin=2em]
\item \textbf{不可逆性}：由于 PID 基于 HMAC-SM3 生成，而 HMAC 的输出在不知道密钥 $K_{\mathrm{master}}$ 的情况下无法逆推出输入 $\mathit{userID} \parallel r$，因此攻击者即使获取了 PID，也无法直接恢复用户的真实身份。该性质基于 SM3 的抗原像性和 HMAC 的 EUF-CMA 安全性。
\item \textbf{不可链接性}：由于每次 PID 生成引入了不同的随机数 $r$，即使攻击者获取了同一用户的多个 PID，也无法判断这些 PID 是否指向同一用户。该性质基于 HMAC 输出的伪随机性。
\end{itemize}

PID 的作用不是取代用户登录身份，而是在订单通信过程中隐藏真实身份。用户在系统中仍然需要通过用户名、口令和设备识别等方式完成正常登录；登录成功后，系统在内部使用 PID 与订单进行绑定。骑手端只需要知道订单编号、订单状态和必要配送信息，不应看到用户真实身份，也不应直接看到 PID。管理员端默认也不直接展示 PID，而是展示经过摘要化处理的 PID 哈希值。

因此，FoodShield 中 PID 的设计原则可以概括为三点。第一，PID 是系统内部匿名标识，不作为普通前端页面中的公开身份展示。第二，PID 与真实身份之间的映射关系只能由平台服务器在受控条件下查询。第三，当发生投诉、恶意骚扰、纠纷处理或审计需求时，系统可以通过条件溯源流程恢复必要身份信息，从而实现隐私保护与责任追踪之间的平衡。
